\begin{center}
    \section*{\fontsize{20}{20}\selectfont Chapter 1}
\end{center}
\vspace{10mm}
\section{Introduction}

Accurate brain tumour segmentation from MRI scans is essential for diagnosis, treatment planning, and longitudinal monitoring of therapeutic response. The BraTS (Brain Tumour Segmentation) challenge has established itself as the standard benchmark for evaluating automated segmentation algorithms on multi-parametric MRI data. Binary segmentation — detecting tumour versus non-tumour tissue — is a fundamental prerequisite for screening workflows, rapid clinical assessment, and initial diagnostics, forming the basis upon which multi-class sub-region delineation builds.

Most current methods rely on volumetric CNNs, particularly 3D U-Net architectures. These approaches achieve strong performance but impose substantial computational demands: 60--80 million parameters and 8+ GB of GPU memory. These requirements create significant barriers to deployment in resource-limited clinical environments, on edge devices, and in settings where rapid inference is essential. There is therefore a clear need for approaches that maintain competitive accuracy while substantially reducing computational cost.

Graph Neural Networks (GNNs) offer a compelling alternative paradigm. Rather than processing entire 3D volumes, GNNs represent images as graphs — superpixels become nodes and spatial relationships become edges — capturing local texture and global structure while operating on substantially smaller data representations. This is a natural fit for medical images of brain tumours, which typically occupy only 5--10\% of total brain volume, creating inherent sparsity that graph-based representations can exploit directly.


\subsection{Problem Statement}

Automated brain tumour segmentation is clinically important, yet current deep learning methods face a persistent efficiency-accuracy trade-off. Volumetric CNNs achieve high accuracy but require substantial computational resources — 68M+ parameters, over 8 GB of GPU memory, and 80+ ms per patient inference time — rendering them impractical for most real-world clinical settings. Lightweight alternatives, conversely, typically sacrifice too much accuracy to be clinically useful.

The central research question is therefore: \textit{Can a graph-based neural network achieve competitive segmentation accuracy ($>$90\% Dice) while remaining efficient enough for deployment in resource-constrained clinical environments?}

Addressing this question requires resolving several technical challenges. First, 3D medical images must be converted into meaningful graph representations without losing diagnostically relevant features. Second, the neural architecture must be compact yet sufficiently expressive for reliable tumour detection. Third, validation must ensure zero data leakage through rigorous patient-level experimental protocols. Finally, efficiency gains must be quantified through systematic benchmarking against established baselines.

\subsection{Problem Background}

The BraTS (Brain Tumor Segmentation) challenge, established in 2012, has become the de facto benchmark for evaluating brain tumor segmentation algorithms. The BraTS 2021 dataset comprises 1,251 multi-institutional glioma cases, each with four MRI modalities (T1, T1 contrast-enhanced, T2, FLAIR) and expert-annotated ground truth segmentations. The challenge focuses on three nested tumor sub-regions: enhancing tumor (ET), tumor core (TC), and whole tumor (WT).

Traditional approaches relied on handcrafted features and classical machine learning methods (SVM, Random Forests), achieving Dice scores around 70-80\%. The deep learning revolution, particularly 3D U-Net architectures introduced in 2015, improved performance to 85-90\% Dice for state-of-the-art systems. Recent winners of BraTS challenges typically employ ensemble methods combining multiple volumetric CNNs, achieving 88-92\% Dice but at the cost of even greater computational requirements.

Graph Neural Networks emerged as a promising alternative in the broader computer vision community, with successful applications in social network analysis, molecular property prediction, and recommendation systems. However, their application to medical imaging remains limited, with only a handful of preliminary studies showing feasibility but often lacking rigorous validation or efficiency analysis. Critical questions remain about optimal graph construction strategies, feature engineering for medical data, and whether GNNs can match CNN performance while delivering on their efficiency promises.

\subsection{Research Objectives}

The primary objective of this research is to develop and rigorously validate an efficient graph-based neural network approach for brain tumor segmentation that achieves competitive accuracy while providing substantial computational advantages. Specific objectives include:

 \begin{itemize}
    \item Develop a robust graph construction pipeline that converts 3D MRI volumes into 2D graph representations, preserving critical tumor features while reducing computational complexity through superpixel-based node generation
    
    \item Design and optimize a graph neural network architecture specifically for binary brain tumor segmentation, balancing model capacity with parameter efficiency
    
    \item Conduct rigorous 5-fold cross-validation with patient-level stratified splits to ensure generalization capability and statistical validity, implementing comprehensive data integrity checks
    
    \item Perform extensive efficiency analysis comparing inference speed, memory consumption, and model size against established U-Net baseline to quantify practical deployment advantages
    
    \item Execute systematic ablation studies to validate architectural design choices (network depth, hidden dimensions, batch size) and identify optimal configurations
    
    \item Implement ensemble methods to explore upper performance bounds and demonstrate robustness through model averaging strategies
    
    \item Ensure complete reproducibility through deterministic training protocols, seed management, and comprehensive auditing procedures (targeting 15/15 integrity checks)
    
    \item Generate qualitative visualizations and quantitative analyses suitable for clinical interpretation and publication-quality reporting
 \end{itemize}

% Placeholder for figure reference
 
\subsection{Motivations}

Several factors motivated this research direction:

\textbf{Clinical Deployment Challenges:} Despite impressive accuracy, current deep learning models for brain tumour segmentation are difficult to deploy in practice. Hardware requirements, slow inference, and the complexity of integrating with existing hospital Picture Archiving and Communication Systems (PACS) create substantial operational barriers.

\textbf{Resource-Constrained Settings:} Many healthcare facilities, particularly in developing regions, lack access to the high-end GPUs required by volumetric CNNs. This creates a diagnostic equity problem in which AI-assisted analysis is available only to well-funded institutions.

\textbf{Real-Time Applications:} Emerging clinical applications — including intra-operative guidance, emergency screening, and telemedicine — require near-instant segmentation results. Current volumetric methods cannot consistently meet these latency requirements.

\textbf{Structural Promise of Graphs:} Tumours are spatially sparse, typically occupying only 5--10\% of total brain volume. Graph representations that exploit this sparsity have the potential to achieve meaningful efficiency gains relative to dense volumetric processing.

\textbf{Data Integrity Concerns:} Recent literature has identified concerning instances of data leakage in medical imaging studies, where patient information inadvertently crosses training and evaluation boundaries. This motivated the implementation of comprehensive auditing procedures, resulting in 15 out of 15 integrity checks passing.

\textbf{Sustainable and Efficient AI:} The broader AI research community has increasingly prioritised efficiency alongside accuracy. This work investigates whether medical imaging can benefit from similar efficiency improvements without sacrificing clinical utility.


\subsection{Significance of the Research}

This research makes significant contributions to both the medical imaging and machine learning communities:

\textbf{Clinical Impact:} By achieving 92.92\% Dice with 6.9$\times$ faster inference and 156$\times$ parameter reduction, our approach enables deployment in resource-constrained clinical settings, emergency departments, and mobile diagnostic units where current methods are impractical.

\textbf{Methodological Advancement:} We demonstrate that graph-based approaches can match volumetric CNN performance on medical segmentation tasks, challenging the prevailing assumption that 3D convolutions are necessary for medical image analysis.

\textbf{Efficiency Paradigm:} Our work contributes to the growing body of evidence that efficiency and accuracy need not be mutually exclusive, particularly relevant as healthcare AI scales to billions of global users.

\textbf{Validation Rigor:} The comprehensive integrity validation framework (15/15 checks including feature dimension verification, patient leakage detection, label validity, and normalization checks) sets a higher standard for medical AI research reproducibility.

\textbf{Ablation Insights:} Discovery of batch size sensitivity in medical imaging (batch 32 optimal, degrading to 83\% at batch 64) and architectural validation (5 layers = 6 layers, no benefit from added depth) provides actionable guidance for future medical imaging GNN designs.

\textbf{Open Science:} Complete pipeline documentation, reproducible code with deterministic training (seed 42), and detailed methodology enable other researchers to build upon this work and validate our findings.

\subsection{Research Contribution}

This thesis makes the following specific contributions:

\begin{itemize}

    \item \textbf{Novel Graph Construction Pipeline:} A superpixel-based graph generation method using SLIC algorithm with 200 adaptive slices per patient and 15 engineered features (intensity statistics, spatial information, shape/texture descriptors) that captures tumor characteristics without ground-truth leakage.
    
    \item \textbf{Optimized GNN Architecture:} A 5-layer GraphSAGE model with 256 hidden dimensions and 439,041 parameters, validated through systematic ablation showing no benefit from deeper networks (6 layers: 84.00\% vs. 5 layers: 84.03\%).
    
    \item \textbf{Strong Performance with Efficiency:} 90.39\% $\pm$ 0.69\% Dice on 5-fold cross-validation and 92.92\% with ensemble, while achieving 6.9$\times$ faster inference (12.7ms vs. 87.8ms) and 156$\times$ parameter reduction (439K vs. 68M) compared to 3D U-Net.
    
    \item \textbf{Comprehensive Validation Framework:} Patient-level stratified 5-fold splits with zero data leakage, verified through paranoid auditing (15/15 integrity checks: feature dimensions, patient overlap, label validity, normalization, seed consistency, batch configuration).
    
    \item \textbf{Batch Size Discovery:} Empirical finding that smaller batches (32) significantly outperform larger batches (64: 83\%, 48: 86\%) in medical imaging GNN tasks, providing actionable insights for future research.
    
    \item \textbf{Complete Reproducibility Package:} Deterministic training protocols (seed 42), comprehensive pipeline documentation (8-phase guide), and validated code enabling full experimental reproduction.
    
    \item \textbf{Clinical Deployment Roadmap:} Demonstration that real-time inference (12.7ms), low memory footprint (2.1 GB), and small model size (439K parameters) enable deployment on edge devices, mobile platforms, and resource-constrained clinical environments.
\end{itemize}


\subsection{Complex Engineering Analysis}

This research addresses complex engineering challenges requiring depth of knowledge across graph neural networks, medical image processing, and clinical validation. The work navigates conflicting requirements between accuracy (target $>$90\% Dice) and efficiency (target $<$50ms inference, $<$1M parameters), conducts extensive analysis through stratified 5-fold cross-validation with 15 integrity checks, and tackles domain-specific issues including MRI protocol variability, extreme class imbalance (tumors occupy 5-10\% of brain volume), and zero data leakage requirements critical for valid medical AI.

The project integrates diverse resources (BraTS 2021 dataset, PyTorch Geometric, SimpleITK, consumer-grade RTX 2060 6GB GPU) and requires interdisciplinary interaction bridging computer science, medical imaging, and clinical practice. Novel contributions include the first comprehensive GNN efficiency study for brain segmentation, discovery of batch size sensitivity in medical imaging GNN tasks, and demonstration that graph-based methods achieve CNN-competitive accuracy (92.92\%) with 6.9$\times$ speedup and 156$\times$ parameter reduction. These efficiency gains enable deployment in resource-constrained settings (developing countries, rural clinics), directly impacting healthcare accessibility and contributing to sustainable AI through reduced energy consumption at scale.

\subsubsection{Complex Engineering Problems (P1-P7)}

This work addresses seven categories of complex engineering problems as defined by accreditation standards (ABET, BAETE):

\textbf{P1: Depth of Knowledge Required} - Advanced understanding across graph neural networks (GraphSAGE, GAT, message-passing), medical image processing (SLIC superpixels, 3D MRI preprocessing, multi-modal intensity analysis), deep learning optimization (Adam, batch normalization, learning rate scheduling), and software engineering (PyTorch Geometric, CUDA optimization, NIfTI/DICOM standards, version control).

\textbf{P2: Range of Conflicting Requirements} - Navigation of competing objectives: accuracy vs. speed ($>$90\% Dice while maintaining $<$50ms inference), expressiveness vs. efficiency (tumor detection within $<$1M parameters), feature richness vs. construction speed (15 engineered features without excessive overhead), and validation rigor vs. iteration speed (5-fold cross-validation with integrity checks).

\textbf{P3: Depth of Analysis Required} - Extensive experimental validation through stratified 5-fold cross-validation with patient-level splits (900 train, 100 validation, 251 test per fold), systematic ablation studies (5 vs. 6 layers, batch sizes 32/48/64), hyperparameter tuning across 20+ configurations, comprehensive integrity auditing (15 validation checks ensuring zero data leakage), and multi-metric evaluation (Dice, sensitivity, specificity, precision, inference time, memory, parameters).

\textbf{P4: Familiarity with Domain-Specific Issues} - Addressing critical medical imaging challenges including protocol variability (1,251 multi-institutional MRI scans with varying acquisition parameters), image artifacts (noise, motion, intensity inhomogeneities), extreme class imbalance (tumors occupy 5-10\% of brain volume requiring weighted loss), privacy regulations (HIPAA/GDPR compliance), clinical consequences (false negatives delay treatment; false positives cause unnecessary interventions), and zero ground-truth leakage requirements.

\textbf{P5: Extent of Applicable Standards} - Compliance with medical imaging formats (NIfTI for research, DICOM for clinical deployment), data privacy regulations (HIPAA, GDPR), BraTS challenge evaluation protocols, reproducibility standards (seed management, deterministic CUDA operations, dependency versioning), and novel integrity auditing procedures where established benchmarks for GNN medical segmentation did not exist.

\textbf{P6: Level of Uncertainty} - High unpredictability from tumor heterogeneity (extreme variation in glioma appearance, shape, size, intensity, texture), scanner variability (different magnetic field strengths, acquisition protocols, vendor hardware), patient factors (motion during scanning, anatomical variations, pathological diversity), architectural uncertainty (optimal graph construction strategy unclear), and generalization risk (training on 1,251 cases may not cover full clinical diversity).

\textbf{P7: Component Interdependence} - Tightly coupled system requiring orchestration: preprocessing impacts feature quality, graph construction defines model input structure, architecture depends on graph properties, training protocol affects convergence, ensemble strategy requires sufficient model diversity, and clinical integration necessitates standardized outputs for PACS systems and visualization tools.

\subsubsection{Complex Engineering Activities (A1-A5)}

The research demonstrates five categories of complex engineering activities:

\textbf{A1: Diversity of Resources} - Integration of diverse computational and domain resources including BraTS 2021 dataset (1,251 multi-parametric MRI volumes with expert annotations), software libraries (PyTorch Geometric for GNN, SimpleITK for medical imaging, scikit-image for superpixels), consumer-grade hardware (NVIDIA RTX 2060 6GB VRAM demonstrating accessibility), neuroimaging literature (glioma biology, MRI physics), and ML validation frameworks (cross-validation, ablation, ensembles).

\textbf{A2: Level of Interdisciplinary Interaction} - Continuous collaboration across fields: understanding clinical requirements through radiologist consultation for accuracy and speed needs, translating medical knowledge into graph features and network design, ensuring data integrity through validation protocols preventing medical AI pitfalls, interpreting predictions in clinical context (sensitivity to false negatives vs. positives), balancing computer science reproducibility practices with medical privacy and validation standards, and communicating efficiency-accuracy trade-offs to diverse stakeholders.

\textbf{A3: Innovation and Novelty} - Novel technical contributions including first optimized superpixel-GNN integration for 3D medical images with 15 multi-modal features, comprehensive GNN efficiency vs. accuracy trade-off analysis for brain segmentation, characterization of batch size sensitivity in medical imaging GNN tasks (batch 32 optimal, 64 degrades), systematic ablation showing 5-layer architectural sweet spot (6 layers provide no benefit), 15 comprehensive integrity checks addressing reproducibility crisis in medical AI, and demonstration of CNN-competitive accuracy (92.92\%) with 6.9$\times$ speedup and 156$\times$ parameter reduction.

\textbf{A4: Societal and Environmental Consequences} - Significant implications for healthcare access (enables deployment in developing countries, rural clinics, mobile diagnostic units where volumetric CNNs impractical), clinical efficiency (6.9$\times$ faster inference reduces patient waiting, enables real-time intra-operative guidance), energy sustainability (156$\times$ smaller models reduce carbon footprint at scale), clinical applications (improved segmentation assists diagnosis, treatment planning for surgery and radiation, therapeutic monitoring), error consequences (direct clinical impact requiring careful validation and uncertainty quantification), privacy protection (ethical patient data handling), and global impact (democratizing AI-assisted medical diagnosis worldwide).

\textbf{A5: Knowledge Application and Self-Directed Learning} - Autonomous learning beyond standard curriculum including MRI physics and T1/T2/FLAIR contrasts, glioma pathology and tumor growth patterns, NIfTI structures and skull-stripping algorithms, graph theory message passing and attention mechanisms, patient-level splitting and integrity checks for medical AI validation, GPU memory management and mixed-precision training, and bridging computer science, medical imaging, and clinical practice through effective communication.

 
\subsection{Thesis/Project Organization}

This thesis is organized into six chapters that provide a comprehensive and systematic presentation of the research methodology, experimental results, and contributions:

\textbf{Chapter 1: Introduction} presents the research context, motivation, and scope. It introduces brain tumor segmentation challenges, outlines the efficiency-accuracy dilemma in current deep learning methods, and explains why graph neural networks offer a promising alternative. The chapter defines specific research objectives, describes the significance of the work for clinical deployment, and details the complex engineering aspects of the project.

\textbf{Chapter 2: Literature Review} surveys existing approaches in medical image segmentation, focusing on CNN-based methods (U-Net, V-Net, attention mechanisms) and emerging graph-based techniques. It critically analyzes state-of-the-art performance on BraTS benchmarks, reviews efficiency optimization strategies, discusses data validation concerns in medical AI research, and identifies gaps that motivate the current work.

\textbf{Chapter 3: Methodology} provides detailed technical descriptions of the proposed approach. It documents the BraTS 2021 dataset characteristics, presents the complete pipeline from MRI preprocessing to graph construction, explains the GNN architecture design and training procedures, describes the patient-level stratified cross-validation strategy, and outlines integrity validation protocols that ensure zero data leakage.

\textbf{Chapter 4: Results and Analysis} presents comprehensive experimental findings. It reports 5-fold cross-validation performance (90.39\% $\pm$ 0.69\% Dice), ensemble results (92.92\% Dice), efficiency benchmarking against U-Net baselines (6.9$\times$ speedup, 156$\times$ parameter reduction), systematic ablation studies validating architectural choices, batch size sensitivity analysis, and qualitative visualizations of segmentation outputs.

\textbf{Chapter 5: Discussion} interprets the results in broader context. It compares performance against state-of-the-art methods, analyzes the implications of efficiency gains for clinical deployment, discusses the significance of architectural findings (5 vs. 6 layers, batch size effects), acknowledges limitations (preprocessing overhead, feature engineering requirements), addresses potential clinical translation barriers, and suggests directions for future work.

\textbf{Chapter 6: Conclusion} synthesizes the key findings and contributions. It summarizes how the research addresses the stated objectives, highlights the main contributions to graph-based medical imaging, emphasizes the practical impact of efficiency improvements, and reflects on lessons learned about balancing accuracy with computational constraints in medical AI applications.


\vspace{5mm}

This chapter has established the foundation for a comprehensive investigation into efficient graph-based brain tumour segmentation. The critical gap between high-performing but computationally expensive volumetric CNNs and the clinical need for efficient, deployable diagnostic tools has been identified and motivated. The specific objectives outlined provide a roadmap for rigorous validation — from graph construction through to ensemble methods — with emphasis on reproducibility and data integrity. By framing this work within complex engineering problem-solving contexts, the interdisciplinary nature of medical AI research bridging computer science, medical imaging, and clinical practice has been underscored. The subsequent chapters demonstrate how graph neural networks can achieve competitive accuracy (92.92\% Dice) while delivering substantial efficiency improvements (6.9$\times$ inference speedup, 156$\times$ parameter reduction) suitable for real-world clinical deployment.